% ==============================================================================
% SECTION 5: CONCLUSIONS
% ==============================================================================

\section{Conclusiones}\label{sec:conclusions}

\subsection{El cambio de paradigma: de LLMs a sistemas agénticos}
\label{subsec:paradigm-shift}

La revisión de la literatura permite concluir que los modelos de lenguaje por sí
solos son insuficientes para las demandas de la bioinformática debido a su
naturaleza probabilística. Sin embargo, la transición hacia sistemas agénticos,
que integran razonamiento, uso de herramientas externas y ciclos de
retroalimentación, ha demostrado ser un avance disruptivo.

Estos sistemas no solo ``escriben'' código, sino que actúan como agentes
ejecutores capaces de iterar hasta alcanzar soluciones técnicas válidas, como se
observó en la automatización del modelado farmacocinético y el análisis de
\textit{single-cell omics}.

\subsection{La verificabilidad como requisito científico}
\label{subsec:verifiability}

Un hallazgo crítico es que la autonomía no debe comprometer la integridad
científica. La implementación de estrategias de \textit{grounding} y
verificación a nivel de sentencia, como en el caso de Alvessa, es esencial para
transformar las alucinaciones de los modelos en resultados trazables.

La integración de herramientas de observabilidad como LangSmith se perfila no
como un lujo técnico, sino como una necesidad metodológica para garantizar la
auditoría y la reproducibilidad de los descubrimientos científicos generados por
IA.

\subsection{Perspectivas futuras y estandarización}
\label{subsec:future-standardization}

El futuro de la automatización en bioinformática depende de la creación de un
ecosistema interoperable. La adopción de protocolos como el \textit{Model
Context Protocol} (MCP) y el desarrollo de benchmarks especializados, por
ejemplo GenomeArena, serán los catalizadores que permitan pasar de prototipos
aislados a infraestructuras de investigación robustas.

A medida que se reduzcan los costes computacionales y mejore la precisión de los
agentes, estos sistemas se convertirán en colaboradores esenciales en el
laboratorio, permitiendo que el investigador humano se desplace de las tareas de
ejecución técnica hacia la formulación de hipótesis de alto nivel y la toma de
decisiones estratégicas.

En conclusión, los sistemas basados en agentes no solo automatizan tareas
repetitivas, sino que democratizan el acceso a flujos de trabajo bioinformáticos
complejos, siempre que su despliegue esté regido por marcos de gobernanza
transparentes, observabilidad estricta y una validación humana continua.
